\uuid{GHCq}
\exo7id{3880}
\titre{exo7 3880}
\auteur{quercia}
\organisation{exo7}
\datecreate{2010-03-11}
\isIndication{false}
\isCorrection{false}
\chapitre{Continuité, limite et étude de fonctions réelles}
\sousChapitre{Etude de fonctions}
\cpge{mpsi-mp2i.fonctions-valeurs-reelles-complexes.a, pcsi.techniques-calcul-fonctions.a, ptsi.techniques-calcul-fonctions.a, mpsi-mp2i.limites-continuite.a, pcsi.limites-continuite.a, ptsi.limites-continuite.a, mpsi-mp2i.limites-continuite.c, pcsi.limites-continuite.c, ptsi.limites-continuite.c, bcpst1.fonctions-reelles-usuelles.a, bcpst1.limites-continuite.a, bcpst1.limites-continuite.c}
\module{Analyse}
\niveau{L1}
\difficulte{}

\contenu{
\texte{
Soit $f : {]0,1]} \to {\R}$ uniformément continue.
Pour $x \in {]0,1]}$ on pose :
$\begin{cases} g(x) = \sup\bigl(f(]0,x])\bigr) \cr
         h(x) = \inf\bigl(f(]0,x])\bigr).\cr\end{cases}$
}
\begin{enumerate}
    \item \question{Montrer que $g$ et $h$ sont monotones.

On note $\ell = \lim_{x\to0^+} g(x)$ et $m = \lim_{x\to0^+} h(x)$.}
    \item \question{En utilisant la continuité uniforme de $f$, montrer que $\ell = m$.}
    \item \question{En déduire que $f(x) \to \ell$, si $x\to0^+$.}
\end{enumerate}
}
